\documentclass[12pt]{article}
\usepackage[utf8]{inputenc}

\usepackage{../mystyle}
\usepackage{parskip}
\usepackage{float}
\usepackage[normalem]{ulem}

\usepackage{blkarray}

% Comic Sans
% \usepackage{sansmathfonts}
% \usepackage{fontspec}
% \setmainfont{Comic Sans MS}

% \usepackage{fontspec}
% \setmainfont{Wingdings-Regular.ttf}

% \usepackage{libertine}

% fancy header
\usepackage{fancyhdr}
\pagestyle{fancy}
% \rhead{Joseph Sullivan}

\title{Tilting}
\author{Joseph Sullivan}
\date{March 2025}



\begin{document}

\maketitle
\tableofcontents

\section{Introduction}

This set of notes tries to explain a certain correspondence between coherent sheaves and quiver representations.

To understand this story, we'll need to do understand
\begin{itemize}
    \item endomorphisms of certain nice generator sheaves, called tilting sheaves
    \item how to compute with quiver representations.
\end{itemize}

These notes will assume some working familiarity with triangulated categories---the reference \cite{calThesis} has a nice treatment, also see my other notes.

These notes are also not really intended to be read from start to finish---skip the quiver section and go back as needed.

\section{Quiver Algebras}\label{quiver-section}

We're going to need to understand a bit of the representation theory of quivers. This is more or less copied from \cite{quiver}, except we switched the convention of path composition and from right modules to left modules.

\subsection{Definitions}

\begin{defn}
    A \textbf{quiver} $Q=(Q_0,Q_1,s,t)$ is a quadruple consisting of sets $Q_0$ called the \textbf{vertices}, $Q_1$ called the \textbf{arrows}, and maps $s,t:Q_1\to Q_0$ which associate to each arrow $\alpha\in Q_1$ its \textbf{source} $s(\alpha)\in Q_0$ and its target $t(\alpha)\in Q_0$.

    We will assume all quivers are finite and connected (the underlying the undirected graph is connected).
\end{defn}

\begin{defn}
    A \textbf{path} in $Q$ is a sequence
    \[(a \mid \alpha_1,\alpha_2,\dots,\alpha_\ell \mid b)\]
    where $\alpha_k \in Q_1$ are arrows with $t(\alpha_k)=s(\alpha_{k+1})$ (unless $k=0,\ell$, in which case $a=s(\alpha_1)$ and $t(\alpha_\ell)=b$). We allow $\ell=0$ when $a=b$ for a trivial path
    \[\eps_a = (a \| a).\]
\end{defn}

\begin{defn}
    Let $k$ be a field and $Q$ a quiver, the \textbf{quiver algebra} (or \textbf{path algebra}) $kQ$ is the $K$-algebra with underlying vector space
    \[kQ = \bigoplus k\cdot (a\mid \alpha_1,\dots,\alpha_\ell | b)\]
    the free $k$-space on the set of paths in $Q$, and multiplication defined on the generators by
    \[(c \mid \beta_1,\dots,\beta_k \mid d) \cdot (a \mid \alpha_1,\dots,\alpha_\ell \mid b) = \begin{cases}
        (a \mid \alpha_1,\dots,\alpha_\ell,\beta_1,\dots,\beta_k \mid d) & b=d\\
        0 & \text{otherwise},
    \end{cases}\]
    (i.e., composition of paths), and the multiplication is extended to all of $kQ$ by distributivity.

    Denote $kQ_m$ the $k$-vector subspace of $kQ$ spanned by paths of length $m$. This makes $kQ$ into a graded $k$-algebra.
\end{defn}

\begin{rmk}
    The associative $k$-algebra $kQ$ is unital. We claim that $\sum_{a\in Q_0} \eps_a$ is a unit for multiplication, since for a given path $\gamma$, it will end at some $t(\gamma)$, and then
    \[\eps_a \cdot \gamma =\begin{cases}
        \gamma & a = t(\gamma)\\
        0 & \text{otherwise},
    \end{cases}\]
    so indeed $(\sum \eps_a)\cdot \gamma = \gamma$. The same argument shows $\gamma \cdot (\sum \eps_a) = \gamma$, so we have
    \[\sum_{a\in Q_0} \eps_a = 1.\]
\end{rmk}

\begin{defn}
    Let $Q$ be a quiver and $kQ$ its quiver algebra. Define the \textbf{arrow ideal} of $KQ$ to be $R_Q \defeq \<Q_1\>_{kQ}$ the (two sided) ideal generated by the arrows (paths of length 1). In this case,
    \[R_Q = \bigoplus_{m \geq 1} kQ_m,\]
    and more generally
    \[R_Q^\ell = \bigoplus_{\ell \geq m} kQ_\ell.\]
\end{defn}

\begin{defn}
    A two-sided ideal $I\leq kQ$ is called \textbf{admissible} if there exists $m\geq 2$ such that
    \[R^m_Q \subseteq I \subseteq R^2_Q,\]
    i.e., $I$ is generated by formal sums of paths of length at least $2$, and any long enough path is in $I$.

    We call a pair $(Q,I)$ a \textbf{bound quiver}, and the $k$-algebra $kQ/I$ the \textbf{bound quiver algebra}.

    Every admissible ideal turns out to be generated by finitely many \textbf{relations}, i.e., a $k$-linear sum of paths (of length $\geq 2$) with the same source and target.
\end{defn}

\begin{defn}
    Let $(Q,I)$ be a bound quiver. A \textbf{bound quiver representation} is an assignment $(M_a,\phi_\alpha)$ of $k$-vector spaces $M_a$ for $a\in Q_0$ and $k$-linear maps $\phi_\alpha: M_{s(\alpha)} \to M_{t(\alpha)}$ for $\alpha\in Q_1$ where for each relation $\sum \lambda_i \gamma_i \in I$,
    \[\sum \lambda_i \phi_{\gamma_i} = 0,\]
    where for a path $\gamma = (a\mid \alpha_1\cdots \alpha_\ell \mid b)$,
    \[\phi_\gamma = \phi_{\alpha_\ell} \circ \cdots \circ \phi_{\alpha_1}.\]

    A morphism of bound quiver representations $(M_a,\phi_\alpha) \to (N_a,\psi_\alpha)$ consists of morphisms of vector spaces $M_a\to N_a$ commuting with $\phi_\alpha,\psi_\alpha$. This makes bound quiver representations into an abelian category, denoted
    \[\Rep_k(Q,I).\]
    When all vector spaces are finite dimensional, call the representation a \textbf{finite dimensional} representation, and the full subcategory of these is denoted
    \[\rep_k(Q,I),\]
    which is also abelian.
\end{defn}

\begin{construction}
    Let $(Q,I)$ be a bound quiver, and $(M_a,\phi_\alpha)$ be a bound quiver representation. We get a left $(kQ/I)$-module $M$ as follows. Define the underlying $k$-vector space by
    \[M \defeq \bigoplus_{a\in Q_0} M_a,\]
    and define the $kQ$-action by first defining for $\alpha\in Q_1$
    \[\alpha\cdot \sum_{a\in Q_0 }m_a \defeq \phi_\alpha \cdot m_{s(\alpha)},\]
    and then extending by composition/linearity. One can check that $I$ is in the annihilator of $M$, so this descends to a $kQ/I$ action on $M$.

    Conversely, let $M$ be a left $(kQ/I)$-module. Define a bound quiver representation by setting for $a\in Q_0$
    \[M_a = \eps_a M\]
    and for $\alpha \in Q_1$ with $a=s(\alpha)$, $b=t(\alpha)$, denote $L_\alpha: M\to M$ left multiplication by $\alpha$, and set
    \[\phi_\alpha = L_\alpha|_{M_a} : M_a \to M_b\]
    (where $\alpha\cdot M_a = \alpha \cdot \eps_a M = \alpha\cdot M$ lands in $M_b = \eps_b M$ because $\alpha\cdot M = \eps_b \alpha\cdot M \subseteq \eps_b M$). One can check that $I$ in the annihilator of the $kQ$-action on $M$ will make $(M_a,\phi_\alpha)$ satisfy the relations of $I$.

    One can show that this gives an equivalence of categories
    \[\Rep_k(Q,I) \simeq A\lMod,\]
    which restricts to an equivalence
    \[\rep_k(Q,I) \simeq A\lmod.\]
\end{construction}


\subsection{Projective Resolutions}

To make some explicit calculations of left derived functors out of $D^b(\rep_k(Q,I))$, we will need to learn how to build projective resolutions of quiver representations.

We observe that as left $kQ/I$ modules,
\[kQ/I = \bigoplus_{a\in Q_0} (kQ/I)\eps_a,\]
i.e., every path starts at some unique $a\in Q_0$. This shows that each $(kQ/I)\eps_a$ is a projective object, which we will use to start building our projective resolutions.

\begin{construction}\label{quiver-res}
    We can start any projective resolution as follows. Let $M$ be a left $(kQ/I)$-module, then we first surject onto $M$ by
    \begin{align*}
        \bigoplus_{a\in Q_0} (kQ/I) \eps_a \otimes_k \eps_a M & \xlongrightarrow{\eps} M\\
        \eps_a \otimes m_a &\longmapsto m_a,
    \end{align*}
    (and extended by linearity) which is indeed surjective because
    \[M = 1\cdot M = \left(\sum_{a\in Q_0} \eps_a\right) \cdot M = \bigoplus_{a\in Q_0} \eps_a M.\]
    Now, we need to generate the kernel. We notice some more ``obvious'' elements of the kernel, where for any arrow $\alpha \in Q_1$ and $m_{s(\alpha)} \in \eps_{s(\alpha))} M$,
    \[\alpha \otimes m_{s(\alpha)} - \eps_{t(\alpha)} \otimes (\alpha\cdot m_{s(\alpha)}) \longmapsto \alpha \cdot m_{s(\alpha)} - \alpha \cdot m_{s(\alpha)} = 0.\]
    We can then at least get some of the kernel of $\eps$ by
    \begin{align*}
        \bigoplus_{\alpha\in Q_1} (kQ/I)\eps_{t(\alpha)} \otimes_k \eps_{s(\alpha)} M &\xlongrightarrow{\phi_1} \bigoplus_{a\in Q_0} (kQ/I)\eps_a \otimes_k \eps_a M \tag{$*$}\label{kernel-generation}\\
        \eps_{t(\alpha)} \otimes m_{s(\alpha)} &\longmapsto \alpha \otimes m_{s(\alpha)} - \eps_{t(\alpha)} \otimes (\alpha \cdot m_{s(\alpha)}).
    \end{align*}
    We claim that the image of $\phi_1$ is exactly $\ker \eps$, so that we are successful at least with our first step. We will prove this in the following proposition, which gives us an exact sequence
    \[\bigoplus_{\alpha\in Q_1} (kQ/I)\eps_{t(\alpha)} \otimes_k \eps_{s(\alpha)} M \xlongrightarrow{\phi_1} \bigoplus_{a\in Q_0} (kQ/I)\eps_a \otimes_k \eps_a M \xlongrightarrow{\eps} M \to 0.\]
    When $I=0$ and $Q$ acyclic, we will see that this is left exact, and we are done with our projective resolution, though in general we will need more steps.
\end{construction}

\begin{prop}
    Let $Q$ be acyclic. Then, the image of $\phi_1$ in (\ref{kernel-generation}) is the kernel of $\eps$.
\end{prop}
\begin{proof}
    By construction, $\im \phi_1 \subseteq \ker \eps$, so our task is to show that $\ker \eps/\im \phi_1 = 0$. 
    
    Let $s\in \ker \eps$ be a representative of some element in the quotient. Define the \textit{support} $\supp(s) \subseteq Q_0$ the set of vertices $i\in Q_0$ such that
    \[\eps_i s \neq 0.\]
    We can then rephrase our task as finding $t \in \ker \eps$ with $s\equiv t \mod \ker \im \phi_1$ such that $\supp(t)=\emptyset$. We will prove this by induction---but what to induct on? We produce a filtration of our quiver $Q$ as follows. Since $Q^0 \defeq Q$ is finite and acyclic, it contains a \textit{sink}, i.e., a vertex with no outgoing arrows. Removing this vertex, (and taking the full subquiver) we will still have a finite acylic quiver $Q^1$, and we can repeat, giving us
    \[\emptyset = Q^n \leq \cdots \leq Q^1 \leq Q^0 = Q.\]
    Now, this filtration gives us an index set to induct on---assuming $s$ is supported on $Q^m$, we show $s\equiv t$ where $t$ is supported on $Q^{m+1}$. Denote $\{b\} = Q^{m+1}\setminus Q^m$ the removed sink.

    Now, we write down an expression for $\eps_b s$, and subtract elements from $\im \phi_1 \cap Q^m$ to ``cancel it out'', to get us our desired $t\in Q^{m+1}$.

    We know $\eps_b s$ lives in
    \[\eps_b \cdot \bigoplus_{i\in Q_0} (kQ/I)\eps_i \otimes_k \eps_i M,\]
    so we can write
    \[\eps_b s = \sum_{i\in Q^m_0} \sum_{\substack{\text{paths}\\ \gamma: i\to b}} \lambda_{\gamma} \gamma \otimes m_{\gamma},\]
    where $m_\gamma \in \eps_i M$. Separating out the $i=b$ summand (and since $\eps_b$ is the only path $b\to b$ by acyclicity),
    \[\eps_b s = \eps_b \otimes m_b + \sum_{i\in Q^{m+1}_0} \sum_{\substack{\text{paths}\\ \gamma: i\to b}} \gamma \otimes m_{\gamma}. \tag{$**$}\]
    We are assuming $\eps(s)=0$, so $\eps_b\eps(s)=\eps(\eps_b s) = 0$, which tells us
    \[\eps_b m_b + \sum_{i\in Q^{m+1}_0} \sum_{\substack{\text{paths}\\ \gamma: i\to b}} \gamma m_{\gamma} = 0. \tag{$***$}\]
    Now, using $\im \phi_1 \subseteq \ker \eps$, we can produce elements
    \[\gamma \otimes m_{\gamma} - \eps_b \otimes \gamma\cdot m_{\gamma},\]
    which are also in $Q^m$ (by iteratively moving an arrow across from the left argument to the right argument of the tensor). So, subtracting all of these (ranging over all $i\in Q^{m+1}_0$ and all $\gamma:i\to b$) from $s$, we will get in ($**$) just a single tensor
    \[\eps_b \otimes m_b' = \eps_b \otimes (m_b - \sum \gamma\otimes m_\gamma),\]
    so that ($***$) turns into $m_b'=0$, which tells us for
    \[t \defeq s - \sum_{i\in Q^{m+1}_0} \sum_{\substack{\text{paths}\\ \gamma: i\to b}}( \gamma\otimes m_\gamma - \eps_b \otimes \gamma m_{\gamma}) \equiv s \mod{\im \phi_1}\]
    we get as desired $\eps_b t = 0$. So, $t$ is supported on $Q^{m+1}$, completing our induction.
\end{proof}

\begin{prop}
    Let $Q$ be acyclic. When $I=0$, the map $\phi_1$ in (\ref{kernel-generation}) is injective.
\end{prop}
\begin{proof}
    {\color{red} TODO.}
\end{proof}

\section{Tilting Equivalence}

\subsection{Exceptional Collections and Tilting Sheaves}
\begin{defn}
    Let $\cD$ be a $k$-linear triangulated category.
    \begin{itemize}
        \item An object $E\in \cD$ is \textbf{exceptional} if it has no nontrivial Exts, i.e.,
        \[\Hom_{\cD}(E,E[\ell]) = \begin{cases}
            k \cdot \id_E & \text{if $\ell=0$}\\
            0 & \text{otherwise.}
        \end{cases}\]
        \item A sequence of objects
        \[E_1,\dots,E_n\]
        is \textbf{exceptional} if all objects are exceptional, and there are no ``backwards'' Exts, i.e.,
        \[\Hom_\cD(E_i,E_j[\ell])=\begin{cases}
            k & \text{if $\ell=0$, $i=j$}\\
            0 & \text{if $\ell\neq 0$, $i=j$}\\
            0 & \text{if $i>j$.}
        \end{cases}\]
        \item A sequence is \textbf{full} if it generates $\cD$ (i.e., any strictly full triangulated subcategory containing the sequence is equal to $\cD$).
        \item A sequence is \textbf{strong} if in addition to being exceptional, the remaining (forward) Exts are all concentrated in degree 0, i.e.,
        \[\Hom_\cD(E_i,E_j[\ell]) = \begin{cases}
            k & \text{if $\ell=0,i=j$}\\
            0 & \text{if $\ell\neq 0$.}
        \end{cases}\]
        (so you are still allowed to have ``forwards'' Homs).
    \end{itemize}
\end{defn}

\begin{eg}
    Consider the following sequence in $D^b(\P^n)$
    \[\cO(-n),\cO(-n+1),\dots,\cO.\]
    Making use of the fact
    \[\Ext^\ell(\cL,\cF) = \Ext^\ell(\cO,\cF\otimes \cL^\vee) = H^\ell(\cF\otimes \cL^\vee),\]
    we can check a few properties of the sequence by standard cohomology computations.
    
    First, the sequence is an exceptional collection, since
    \[\Ext^\ell(\cO,\cO(-i)) = H^\ell(\cO(-i)),\]
    and now,
    \begin{itemize}
        \item if $\ell=0$, then we have $H^0(\cO(-i))$, which is $k$ if $i=0$ and $0$ otherwise,
        \item if $\ell=n$, then we have $H^n(\cO(-i)) = H^0(\cO(i-n-1)) = 0$ for $i\leq n$.
        \item for any other value of $\ell$, we get 0.
    \end{itemize}

    Furthermore, this sequence is also strong, since
    \[\Ext^\ell(\cO,\cO(i)) = H^\ell(\cO(i)),\]
    which is always zero for $\ell\not\in \{0,n\}$, and even when $\ell=n$,
    \[H^n(\cO(i)) = H^0(\cO(-i-n-1)) = 0\]
    for $i\geq 1$ (or even $i\geq -n$, but we cannot add more twists to our collection without losing the adjective ``exceptional'').

    Beilinson proved that this sequence is full---this is somewhat tricky, relying on a certain ``resolution of the diagonal'' in $\P^n\times\P^n$ and some Fourier-Mukai machinery.
\end{eg}

\begin{rmk}
    Exceptional objects are particularly simple. For $E\in \cD$ be exceptional, we claim that the triangulated category generated by $E$ is
    \[\<E\> = \{\bigoplus_i E[i]^{\oplus a_i} \mid \text{all but finitely many $a_i=0$}\}.\]
    This is clearly closed under direct sums and shifts. To see that this is closed under cones, consider a morphism
    \[\bigoplus_i E[i]^{\oplus a_i} \to \bigoplus_i E[i]^{\oplus b_i}.\]
    This is given by some matrix in $\Hom(E[i]^{\oplus a_i}, E[j]^{\oplus b_j})$ (as always for morphisms between direct sums), but it is in fact diagonal because $E$ is exceptional, so the cone over this matrix splits as a direct sum of (shifts of) cones over morphisms
    \[E^{a_i} \to E^{b_i}.\]
    which is just a matrix in $k$. By Gaussian elimination, we can hit both sides with automorphisms (which only changes the cone up to isomorphism) so that the matrix only has diagonal $1$ entries, and so our cone splits as a direct sum of cones over
    \[E \xrightarrow{\id_E} E \qquad \text{or} \qquad E \xrightarrow{0} E,\]
    so we either get $0$ or $E\oplus E[1]$. In any case, the cone will be a direct sum of objects we are including, so indeed $\<E\>$ is given by our description.

    This tells us that $\<E\> \cong D^b(k)$, (a priori objects of $D^b(k)$ look like complexes, but because it has homological dimension $\leq 1$, the complexes split as direct sums of their cohomologies, which gives us the same form as $\<E\>$).
\end{rmk}

Full strong exceptional collections are also very simple---they form a particularly nice generating set for the triangulated category. For a single exceptional object, we got one of the easiest $k$-linear triangulated categories there is, which is $D^b(k)$. Now, for a collection, we'll get something a little more complicated, but still very controllable.

We first pass from exceptional collections to a related object called a \textit{tilting object}, but we need to give another definition related to generating a triangulated category.
\begin{defn}
    We say that a full triangulated subcategory $\cD'$ of a triangulated category $\cD$ is \textbf{saturated}/\textbf{\'epaisse} if it closed under direct summands.

    Write $\<E\>$ for the smallest strictly full saturated triangulated subcategory of $\cD$ containing $E$, and say that $E$ \textbf{classically generates} $\cD$ if $\<E\>=\cD$.
\end{defn}

\begin{warning}
    Sometimes people write $\<S\>$ to mean the extension closure of $S$, i.e., the smallest strictly full subcategory containing $S$ and closed under \textit{extensions}, rather than cones/shifts/direct summands. Always check what someone means.
\end{warning}

\begin{rmk}\label{induction-triangulated-cat}
    For any $E\in \cD$, there is actually an explicit description of $\<E\>$. To get a strictly full saturated triangulated subcategory containing $\<E\>$, one might start with $\{E\}$, and then repeatedly throw in shifts, cones, direct sums, direct summands, etc. The miracle is that if you do this countable many times (unioning the results), you actually get all of $\<E\>$ \cite[Lemma 13.36.2]{stacks-project}. This then gives us the following induction strategy to prove properties of $\<E\>$ \cite[Remark 13.36.7]{stacks-project}.

    \textbf{Induction.} Let $T$ be a property enjoyed by objects of $\cD$ (where we write $T(A)$ to mean $T$ holds for $A$), and suppose
    \begin{enumerate}[(1)]
        \item (base case) $T(E[n])$ for all $n\in \Z$,
        \item if $T(K)$ and $T(L)$, then $T(K\oplus L)$,
        \item if $T(K\oplus L)$ then $T(K)$ and $T(L)$,
        \item if $K\to L \to M \to K[1]$ is a distinguished triangle and $T$ holds for two of them, then $T$ holds for the third,
    \end{enumerate}
    then $T$ holds for all objects in $\<E\>$.
\end{rmk}

\begin{defn}
    Let $X$ be a smooth projective variety over $k$. A sheaf $T\in \Coh(X)$ is called a \textbf{tilting sheaf} if
    \begin{enumerate}
        \item[(T1)] the algebra $A\defeq \End(T)$ has finite global dimension (i.e., there exists $d\in \Z_{\geq 0}$ such that any module admits a projective resolution of length less than $d$).
        \item[(T2)] The modules $\Ext^\ell(T,T)=0$ for all $\ell>0$.
        \item[(T3)] The object classically generates $D^b(X)$.
    \end{enumerate}
\end{defn}

In fact, more is true. An object $X$ of a triangulated category is called \textit{compact} if $\Hom(X,-)$ preserves all coproducts in the triangulated category. In $\Coh(X)$, all coproducts are products, so $\Hom(X,-)$ automatically preserves coproducts, and so a classical generator will actually be a compact generator.

\begin{prop}
    Let $E_1,\dots,E_r \in \cA$ be a full strong exceptional collection for $D^b(\cA)$. Then,
    \[T = \bigoplus_{i=1}^r E_i\]
    is a tilting object.
\end{prop}
\begin{proof}
    (T2) and (T3) are reasonably quick to prove.

    To see (T1), we can see that
    \[A = \End(T)\]
    is equivalent to the $k$-algebra
    \[\begin{bmatrix}
        \Hom(E_1,E_1)=k & \Hom(E_1,E_2) & \cdots & \Hom(E_1,E_n)\\
        0 & \Hom(E_2,E_2)=k & \cdots & \Hom(E_2,E_n)\\
        \vdots & \vdots & \ddots & \vdots\\
        0 & 0 & \cdots & \Hom(E_n,E_n)=k
    \end{bmatrix}\]
    with multiplication given by matrix multiplication/composition of Homs. So, $A$ is equivalent to the quotient algebra $kQ/\<R\>$ for an acyclic bound quiver $(Q,R)$, and the general theory says that these have finite global dimension.
\end{proof}

\begin{rmk}
    Regarding the algebra structure of $A$, we make some comments. The Jacobson radical $J$ of $A$ will be the strictly upper triangular matrices in our suggestive matrix. To see this, set $J$ the direct sum of all Homs which are not $\Hom(E_i,E_i)$ (the strictly upper triangular matrices), and check (via \cite[I.1.4]{quiver}) that $J$ is nilpotent and $A/J$ is a product of copies of $\C$.
    
    First, $J$ is nilpotent because $J^\ell$ only consists of formal sums of morphisms $E_i\to E_{i+\ell}$, so when $\ell\geq r$, there are no more morphisms to sum over.
    
    Next, $A/J$ is isomorphic to
    \[\prod_{i=1}^r \Hom(E_i,E_i) = \prod_{i=-n}^0 \C,\]
    so indeed a product of fields.

    The Jacobson radical is really the tool to find the quiver $Q$. We set the vertices of $Q$ to be a chosen basis for $A/J$, so conveniently we'll choose the basis to be $\id_{E_1},\dots,\id_{E_r}$. Next, we set the arrows from $E_i$ to $E_j$ to be a chosen basis for $\id_{E_j} (J/J^2) \id_{E_i}$. We give a summary of the theory of quiver representations in section \ref{quiver-section}.
\end{rmk}

\begin{eg}
    Let's understand for
    \[A = \End(\bigoplus_{i=-n}^{0} \cO(i))\]
    how to see $A$ as a bound quiver algebra. Our suggestive matrix is the following. 
    \[\begin{bmatrix}
        k\cdot \id_{\cO(\text{-}n)} & \Hom(\cO(\text{-}n),\cO(\text{-}n+1)) & \cdots & \Hom(\cO(\text{-}n),\cO)\\
        0 & k \cdot \id_{\cO(\text{-}n+1)} & \cdots & \Hom(\cO(\text{-}n+1),\cO)\\
        \vdots & \vdots & \ddots & \vdots\\
        0 & 0 & \cdots & k\cdot \id_{\cO}
    \end{bmatrix}= \begin{bmatrix}
        k & k^{\binom{n+1}{1}} & \cdots & k^{\binom{n+n}{n}}\\
        0 & k & \cdots & k^{\binom{n+n-1}{n-1}}\\
        \vdots & \vdots & \ddots & \vdots\\
        0 & 0 & \cdots & k
    \end{bmatrix}\]
    As explained earlier, $J$ is the upper triangular matrices. To compute the arrows from $\cO(i)$ to $\cO(j)$, we compute
    \[\id_{\cO(j)} (J/J^2) \id_{\cO(i)} = \id_{\cO(j)}J \id_{\cO(i)}/ \id_{\cO(j)} J^2 \id_{\cO(i)}.\]
    As a first step, we compute $\id_{\cO(j)} J \id_{\cO(i)}$, which is exactly $\Hom(\cO(i),\cO(j))$. Next, we compute
    \[\id_{\cO(j)} J^2 \id_{\cO(i)},\]
    which will be the subset of $\Hom(\cO(i),\cO(j))$ that factor as a sum of maps which go through some strictly intermediate $\cO(\ell)$. If $j=i+1$, then no strict factorization is possible. If $j>i+1$, then we always have a strict factorization, since a map $\cO(i)\to \cO(j)$ will be given by a degree $j-i$ homogeneous polynomial, which we can write as a sum of monomials, and then peeling one monomial away we get the factorization.

    So,
    \[\id_{\cO(j)} (J/J^2) \id_{\cO(i)} \cong \begin{cases}
        \Hom(\cO(i), \cO(i+1)) & j=i+1\\
        0 & \text{otherwise.}
    \end{cases}\]
    Therefore, we should place $n+1$ arrows between $\cO(i)$ and $\cO(i+1)$ for a chosen basis
    \[\Hom(\cO(i),\cO(i+1)) = \<x_{i,0},\dots,x_{i,n}\>.\]

    So far, we have the quiver $Q$, now we want to compute the relations, so we study the natural map
    \[kQ \to A\]
    given by the universal property of quiver algebras \cite[II.1.8]{quiver}, and we find a generating set of the kernel. We claim that the kernel is generated by terms
    \[R=\{x_{i+1,p}x_{i,q} - x_{i+1,q}x_{i,p} \mid 0\leq p,q \leq n, -n \leq i \leq 0\}.\]
    We can see that all of $R$ is in the kernel, since
    \[x_{p}x_{q} = x_{q}x_{p} \in \Hom(\cO(i),\cO(i+2)).\]
    Now, we check that $\<R\>$ is the entire kernel by seeing that $kQ/\<R\> \to A$ is injective. In fact, both $k$-algebras here are graded by lengths of paths, so we can show the kernel is 0 by showing each graded piece is 0.

    Suppose
    \[\sum \alpha_{p_0,\dots,p_r} x_{i+r,p_r} \cdots x_{i,p_0} \in (kQ/\<R\>)_r\]
    is in the kernel, i.e.,
    \[\sum \alpha_{p_0,\dots,p_r} x_{p_r} \cdots x_{p_0} = 0 \in \Hom(\cO(i),\cO(i+r)) = k[x_0,\dots,x_n]_r,\]
    i.e., it is the zero polynomial. This will imply our original term in $kQ/\<R\>$ is 0, since we can use $R$ to put all the monomials in standard order.

    In conclusion, we get that $A$ is the algebra for the (bound) \textbf{Beilinson $\P^n$ quiver}, i.e., the quiver
    \[\begin{tikzcd}[column sep=1cm]
        \overset{0}{\bullet} \rar[bend left=1cm, shift left=0.5cm,"x_{1,0}"] \rar[bend left=0.5cm, shift left=0.25cm,"x_{1,1}"] \rar[shift left=0cm, phantom, "\vdots"] \rar[bend left=-1cm, shift left=-0.5cm,"x_{1,n}"] & \overset{1}{\bullet} \rar[bend left=1cm, shift left=0.5cm,"x_{2,0}"] \rar[bend left=0.5cm, shift left=0.25cm,"x_{2,1}"] \rar[shift left=0cm, phantom, "\vdots"] \rar[bend left=-1cm, shift left=-0.5cm,"x_{2,n}"] & \overset{2}{\bullet} \rar[phantom,"\cdots"] & \bullet \ar[phantom, "n\,\text{-}1"{font=\scriptsize,above=2pt}] \rar[bend left=1cm, shift left=0.5cm,"x_{n,0}"] \rar[bend left=0.5cm, shift left=0.25cm,"x_{n,1}"] \rar[shift left=0cm, phantom, "\vdots"] \rar[bend left=-1cm, shift left=-0.5cm,"x_{n,n}"] & \overset{n}{\bullet}
    \end{tikzcd}\]
    with the relations
    \[x_{p,i+1}x_{q,i} - x_{q,i+1}x_{p,i} = 0.\]
\end{eg}


\subsection{Main Theorem: LF and RG}
Now, we're getting closer to the main theorem, which tells us how nice a tilting sheaf is---in the same way the triangulated category generated by an exceptional object is just $D^b(k)$, we'll see that the triangulated category generated by a tilting sheaf is $D^b(A)$, where $A$ will actually be a bound quiver algebra (though unfortunately we will have to take the opposite quiver with the opposite relations, since we will get a derived equivalence with right modules over $A$).

First, though, we need a lemma.

\begin{lem}
    Let $F:\cC \to \cD$ be an exact functor, which sends a compact generator $C\in \cC$ to a compact generator $F(C)=D$ of $\cD$, such that the induced map
    \[\End_{\cC}^\bullet(C) \to \End_{\cD}^\bullet(D)\]
    is an isomorphism. Then, $F$ is an equivalence.
\end{lem}
\begin{proof}
    To show at least $F$ is an equivalence, we show it is fully faithful and essentially surjective.
    
    To see fully-faithful, first show that the full subcategory
    \[T_C = \{X\in \cC \mid \text{$\Hom(C,X) \to \Hom(FC,FX)$ is a bijection}\}\]
    is a strictly full saturated triangulated subcategory containing $C[i]$ for all $i\in \Z$. This is true---it is clearly closed under direct sums, direct summands, and shifts. To see it is closed under cones, notice that for a distinguished triangle $X\to Y\to Z \to X[1]$ we get a commutative diagram of (long) exact sequence
    \[\begin{tikzcd}
        \cdots \rar & \Hom(C,X[1]) \rar \dar["F_{C{,}X[1]}"] & \Hom(C,Z) \rar \dar["F_{C{,}Z}"] & \Hom(C,Y) \rar \dar["F_{C{,}Y}"] & \cdots\\
        \cdots \rar & \Hom(FC,FX[1]) \rar & \Hom(FC,FZ) \rar & \Hom(FC,FY) \rar & \cdots
    \end{tikzcd}\]
    all vertical morphisms are isomorphisms except for potentially every 3rd morphism, but then we apply the 5 lemma to get $F_{C,Z}$ is also an isomorphism. So, $T_C$ has to be all of $\cC$.

    The same argument shows that
    \[T_X' = \{Y\in \cC \mid F_{X,Y} \text{ is bijective}\}\]
    is thick (since the previous step told us $C[i]\in T_X'$), so it is also all of $\cC$. This establishes that $F$ is fully faithful.

    Next, $F$ is essentially surjective, because it is exact and has a generator in its image. Thus, $F$ is an equivalence.
\end{proof}

\begin{thm}\label{tilting-quiver}
    Let $T$ be a tilting sheaf on a projective smooth $X$ over $k$, with $A=\End(T)$. Then, the following functors form an equivalence.
    \begin{align*}
        RG= \RHom(T,-): D^b(X) &\longrightarrow D^b(\rmod A)\\
        LF= - \otimes_A^L T: D^b(\rmod A) &\longrightarrow D^b(X)
    \end{align*}
\end{thm}
\begin{proof}
    First, we more carefully define the functors.
    
    We have the left exact functor $G=\Hom(T,-): \QCoh(X) \to k\text{-}\mathrm{Vec}$, which actually factors through $\rMod A$ because we have an action by precomposition.
    
    The category $\QCoh(X)$ has enough injectives, so we can right derive the functor to get
    \[\RHom(T,-): D^+(\QCoh(X)) \to D^+(\rMod A),\]
    and usual arguments tell us that this descends to
    \[RG = \RHom(T,-): D^b(X) \to D^b(\rmod A),\]
    (we have an Ext spectral sequence to ensure that we get bounded complexes because $\Coh(X)$ has finite global dimension, and morphisms between coherent sheaves are finite dimensional to ensure that precomposed with $D^b(X) \cong D^b_{\Coh(X)}(\QCoh(X)) \hookrightarrow D^b(\QCoh(X))$ will with us finitely generated $A$ modules).

    In the other direction, we have the functor
    \[F = - \otimes_A T: \rmod A \to \QCoh(X),\]
    in the following sense.
    
    
    that given $M \in \rmod A$, we get then $\underline{M}$ is in $\rMod \underline{A}$ and $T$ is in $\underline{A} \lMod$, so we can form the $\underline{A}$-balanced tensor product, which is then inherits an $\cO_X$-module structure from $T$ (via left multiplication).

    This functor is right exact, and $\rmod A$ has enough projectives, so we can left derive it to get
    \[- \otimes_A^L T: D^-(\rmod A) \to D^-(\QCoh(X)),\]
    and Serre theorems/finite global dimension of $A$ will give sheaf Tor finite dimensions/vanishing to imply that this functor descends to
    \[LF = - \otimes_A^L T: D^b(\rmod A) \to D^b(X).\]

    Now, we need to show these functors form an equivalence. The assumption T2 for $T$ exactly tells us that $RG(T)=A$, and so we also get a (ring) morphism
    \[\End^\bullet(T) \xrightarrow{RF} \End^\bullet(A),\]
    which is an isomorphism because in nonzero degrees both sides are 0 (by Ext vanishing or because $A$ is projective), and then in degree 0, it sends an endomorphism $f:T\to T$ to left multiplication by $f\in A$, which are exactly the endomorphisms of $A$ as a right $A$-module.
    
    So, our lemma tells us that $RG$ is an equivalence, but a-priori we do not know that $LF$ is a quasi-inverse. To see this, we first notice that $F\dashv G$ form an adjunction before we derived anything. To see this (it seems slightly subtle because it varies slightly from usual tensor-hom in the sheafiness of it), we will spell out the unit and counit, and check the triangle identities. We have $\eta: 1_{\rMod A} \Rightarrow GF$ by
    \[\eta_M: M \to \Hom(T,M\otimes_A T)\]
    where we send $m$ to the morphism $t\mapsto m\otimes t$ (which first is a map to the tensor presheaf, but then we postcompose to get a map to the sheafification). Then, we have $\eps: FG \Rightarrow 1_{\QCoh(X)}$ by
    \[\eps_\cE: \Hom(T,\cE) \otimes_A T \to \cE\]
    where we send $\phi \otimes t$ to $\phi(t)$ (which is a map from the presheaf factoring through the sheafification). Then, one checks that the triangle identities hold by computing the compositions on stalks---where sheafification is not a concern.
\end{proof}


\section{Computations}

\subsection{Computation: LF}

Let $X$ be a projective variety with exceptional collection $E_1,\dots,E_r$, and corresponding tilting sheaf
\[T = \bigoplus_{i=1}^r E_i.\]
We'd like to compute this functor
\[LF = (-) \otimes_A^L T: D^b(\rmod A) \to D^b(X)\]
explicitly. Even more, we can concretely understand $\rmod A$ as $\rep_k(Q^\op,R^\op)$, and we'll ultimately want to understand the entire composition
\[D^b(\rep_k(Q^\op,R^\op)) \to D^b(X).\]

\subsubsection{Setup}
To compute $LF$, as usual, we need to replace either $\underline{M}_\bullet$ or $T$ with a quasi-isomorphic complex of flat $\underline{A}$-modules. To replace $\underline{M}_\bullet$, we make the following observation:

If $M$ is flat in $\rmod A$, then $\underline{M}$ is flat in $\rMod \underline{A}$.

To see this, note that $M$ flat in $\rmod A$ implies $M$ is projective because finitely generated modules over Noetherian rings are finitely presented, and any finitely presented flat module is projective. Therefore, $M$ is a direct summand of a free module $A^n$, so it is still projective in $\rMod A$, so it is flat in $\rMod A$.

Now, to check that $\underline{M}$ is flat, we take a short exact sequence in $\rMod \underline{A}$
\[0 \to \cF \to \cG \to \cH \to 0,\]
and apply $\underline{M} \otimes_{\underline{A}} -$ to get
\[0 \to \underline{M} \otimes_{\underline{A}} \cF \to \underline{M} \otimes_{\underline{A}} \cG \to \underline{M} \otimes_{\underline{A}} \cH \to 0,\]
which is still exact because taking stalks at $x\in X$ we get
\[0 \to M \otimes_A \cF_x \to M\otimes_A \cG_x \to M \otimes_A \cH_x \to 0,\]
which is the application of $M\otimes_A -$ to the short exact sequence which is the stalks, and we know by assumption $M$ is flat in $\rMod A$.

So, the moral is that we can replace a given $M_\bullet$ with a quasi-isomorphic complex of flat/projective objects in $\rmod A$ to do our computation, which is actually a practical thing to do, whereas it's hard to understand how to resolve $T$ using projective/flat objects in $\rMod \underline{A}$.


\subsubsection{Projectives}

As in the setup, we should start by understanding how $LF$ maps projective objects. As a first step, we show the following.
\begin{prop}
    $LF(A) \cong A$.
\end{prop}
\begin{proof}
    Since $A$ is projective, $LF(A)$ is computed as $F(A) = A\otimes_A T$. Now, we have an \textbf{evaluation map}
    \begin{align*}
        A\otimes_A T &\xrightarrow{\eval} T\\
        (A\otimes_A T)(U) &\rightarrow T(U)\\
        \phi \otimes t &\mapsto \phi_U(t)
    \end{align*}
    (which, to be precise, is first defined on the tensor presheaf, and then factors through the tensor sheaf). On stalks, this is the map
    \begin{align*}
        A\otimes_A T_x &\longrightarrow T_x\\
        \phi\otimes t &\longmapsto \phi_x(t),
    \end{align*}
    and this is an isomorphism, since this is exactly the left action on $T_x$. In other words, $\phi\otimes t = 1\otimes \phi_x(t)$, so we have an inverse map $t\longmapsto 1_T \otimes t$.
\end{proof}

Next, we want to see how to compute this functor in terms of quiver representations. We saw that
\[A \cong kQ/\<R\>\]
for some quiver $Q$ and relations $R$, so we get
\[\rmod A \cong kQ^\op/\<R^\op\> \lmod \cong \rep_k(Q^\op,R^\op).\]
We have to come to terms with the ``opposite'' flipping the arrows around, since if we really want to study quiver representations we have to consider left modules (with our conventions).

Recall that a quiver representation is an assignment of vector spaces to vertices and linear maps to arrows satisfying the relations. Now that we've inverted the arrows, we can think of our job as assigning functionals to vertices and pullbacks to arrows. To reinforce this convention, for every arrow $\alpha \in Q$, we will label its opposite arrow $\alpha^* \in Q^\op$.

Now, we try to understand how such a representation of $Q^\op$ is sent over to $D^b(X)$. As in our remark, we should resolve our quiver representation using flat/projective objects, so our task boils down to finding the image of projective quiver representations and the image of morphisms.

Let $1_{E_1},\dots,1_{E_r}$ be the identity maps on each object in our exceptional collection. These are exactly the vertices/length 0 paths in $Q$, and so $1_{E_i}^*$ are the vertices/length 0 paths in $Q^\op$.

The projective objects in $A^\op = kQ^\op/\<R^\op\>$ are then given by
\[P_i = A^\op 1_{E_i},\]
which is exactly the $k$-span of (opposite) paths starting at the vertex $1_{E_i}$. In terms of the endomorphism ring, these are the endomorphisms $T\to T$ with image in $E_i$.

Now, we compute $LF$ on $P_i$, and since $P_i$ is projective, it is computed just by $F$. Since $P_i$ is projective, it is a direct summand of $A$, so the injection $P_i \hookrightarrow A$ is carried to an injection
\[P_i \otimes_A T \hookrightarrow A \otimes_A T,\]
where the right is isomorphic to $T$ by the evaluation map. So, we are really just restricting the evaluation map to the subobject $P_i \otimes_A T$.

\begin{prop}
    $LF(P_i) \cong E_i$.
\end{prop}
\begin{proof}
    We compute the image of $\eval$ restricted to $P_i \otimes_A T$. First, we see that $\eval$ factors through $E_i \leq T$, since
    \[\phi \otimes t \mapsto \phi_U(t) \in E_i\]
    when $\phi$ has image in $E_i$. So, we get a well-defined
    \[P_i \otimes_A T \xrightarrow{\eval} E_i.\]
    This is at least an injection, since it is a restriction (of the codomain) of the composition of two injections.
    
    So, to see that this is an isomorphism, we just show that the stalk maps are surjections. The stalk maps are given by
    \begin{align*}
        P_i \otimes_A T_x &\longrightarrow (E_i)_x\\
        \phi \otimes t &\longmapsto \phi_x(t),
    \end{align*}
    and indeed for $t\in (E_i)_x$ if we take $\pi_i:T\to E_i$ to be projection onto $E_i$, then we have preimage $\pi_i \otimes t$.
\end{proof}

Finally, to aid in our computation, we should describe how $F$ sends morphisms $P_i \xrightarrow{f^*} P_j$ to morphisms $E_i \to E_j$, in the sense of making the diagram commute.
\[\begin{tikzcd}
    P_i \otimes_A T \rar["f^*\otimes 1_T"] \dar & P_j \otimes_A T \dar\\
    E_i \rar & E_j
\end{tikzcd}\]
Note that we continue making the convention of thinking of elements of $A^\op$ modules as functionals, so we denote the morphism $P_i\to P_j$ by $f^*$ ({\color{red} this is maybe a bad convention (as with all my other conventions)}). In any case, such a morphism is determined by its restriction
\[1_{E_i}^* P_i \longrightarrow 1_{E_i}^* P_j,\]
which in terms of endomorphisms of $T$, is
\[k\cdot \{\pi_i\} \longrightarrow k\cdot \{\text{morphisms $E_i \to E_j$}\},\]
just a choice of endomorphism $f^*(\pi_i)$ (where $\pi_i$ is really the same as $1_{E_i}$).

Indeed, if we know this restriction, we get from our commutative diagram for a section $s\in E_i$
\[\begin{tikzcd}
    \pi_i \otimes s \rar[mapsto] & f^*(\pi_i) \otimes s \dar[mapsto]\\
    s \uar[mapsto] \rar[mapsto] & f^*(\pi_i)(s)
\end{tikzcd}\]

In summary, $P_i \xrightarrow{f^*} P_j$ is sent to $f^*(1_{E_i}^*) \in 1_{E_i}^* P_j = \Hom(E_i,E_j)$.

\subsubsection{Examples}

Consider
\begin{itemize}
    \item $X=\P^1$,
    \item exceptional collection $\cO(-1),\cO$
    \item tilting sheaf $T=\cO(-1)\oplus \cO$
    \item endomorphism algebra
    \[A=\End_{\cO_{\P^1}}(T) = \begin{bmatrix}
        \Hom(\cO(-1),\cO(-1)) & \Hom(\cO(-1),\cO)\\
        0 & \Hom(\cO,\cO)
    \end{bmatrix}\]
    corresponding to quiver
    \[Q = \left[\begin{tikzcd}
        \overset{1_{\cO(-1)}}{\bullet} \rar[shift left=0.2cm, bend left=0.2cm, "x"] \rar[shift left=-0.2cm, bend left=-0.2cm, "y"'] & \overset{1_{\cO}}{\bullet}
    \end{tikzcd}\right]\]
    and opposite quiver
    \[Q^\op = \left[\begin{tikzcd}
        \overset{1_{\cO(-1)}^*}{\bullet} & \overset{1_{\cO}^*}{\bullet} \lar[shift left=0.2cm, bend left=0.2cm, "y^*"] \lar[shift left=-0.2cm, bend left=-0.2cm, "x^*"']
    \end{tikzcd}\right]\]
\end{itemize}

Now, let
\[M =\left[\begin{tikzcd}
    \overset{V}{\bullet} & \overset{W}{\bullet} \lar[shift left=0.2cm, bend left=0.2cm, "T^*"] \lar[shift left=-0.2cm, bend left=-0.2cm, "S^*"']
\end{tikzcd}\right]\]
be a representation of $Q^\op$. Then, this has a projective resolution as in Construction \ref{quiver-res} by (labeling the direct summands on the left term)
\[0 \to (W \otimes_k P_{-1}) \oplus (W\otimes_k P_{-1}) \to V \otimes_k P_{-1} \oplus W \otimes_k P_0\to M\to 0,\]
where the first map is determined by the restriction to the $k$-vector space
\[W^{\oplus 2} \cong (W\otimes_k 1_{\cO(-1)}^* P_{-1})^{\oplus 2} \hookrightarrow (W\otimes_k P_{-1})^{\oplus 2},\]
which is given by the matrix
\[\begin{blockarray}{ccc}
& W & W\\
\begin{block}{c[cc]}
    V\otimes_k P_{-1} & S^*(-) \otimes 1_{\cO(-1)}^* & T^*(-) \otimes 1_{\cO(-1)}^*\\
    W\otimes_k P_0 & (-) \otimes x^* & (-) \otimes y^* \\
\end{block}
\end{blockarray}\]
Now, truncating $M$, we get a quasi-isomorphic projective complex, which is sent to
\[0 \to W \otimes_k \cO(-1) \oplus W\otimes_k \cO(-1) \to V\otimes_k \cO(-1) \oplus W\otimes_k \cO \to 0,\]
where the map is given by the following matrix
\[\begin{blockarray}{ccc}
& W\otimes_k \cO(-1) & W\otimes_k \cO(-1)\\
\begin{block}{c[cc]}
    V\otimes_k \cO(-1) & S^*\otimes 1 & T^*\otimes 1 \\
    W\otimes_k \cO & 1_W \otimes x & 1_W \otimes y \\
\end{block}
\end{blockarray}\]
which can be seen by fixing an argument $w\in W$ and studying compositions like
\begin{align*}
    P_{-1} \xrightarrow{\cong} w\otimes_k P_{-1} \to S^*(w) \otimes_k P_{-1} \xrightarrow{\cong} P_{-1}
\end{align*}
where we know how to compute the entire composition by the discussion in the setup.


\subsubsection{Simples}

The equivalence $LF: D^b(\rmod A) \to D^b(\P^n)$ endows $D^b(\P^n)$ with a heart of a t-structure by the image
\[\rmod A \hookrightarrow D^b(\rmod A) \to D^b(\P^n).\]
To try to characterize this heart, we'll find all the simples of $\rmod A$, compute their images in $D^b(\P^n)$, and then our heart will be the extension closure of those simples.

This is done more efficiently in the next section, the duality of simples and projectives. Nonetheless we sketch the computation for projective space.

The simple quiver representations are those assigning $k$ to a single vertex, and $0$ to the remaining vertices (and arrows). To compute the image of this representation under $LF$, we need to write a projective resolution of the simple representation, which we claim looks like a twisted truncated from the left Koszul complex. Roughly this is because to surject onto the simple $S_i$ concentrated at $\cO(i)$, we will use one copy of $P_i = A^\op 1_{\cO(i)}$, and this surjection has a kernel $K$ generated by $1_{\cO(i-1)}K$ which is dimension $n+1$ (for the $n+1$ arrows from $\cO(i-1)\to \cO(i)$), and then we will surject onto the kernel using $P_{i-1}$, so our resolution looks like
\[\cdots \to P_{i-1} \otimes k^{n+1} \to P_i \to S_i \to 0,\]
and the kernel $K$ of $P_{i-1} \otimes k^{n+1} \to P_i$ will be generated by $1_{\cO(i-2)} K$, which has dimension 
\[(n+1)^2 - (\binom{n+1}{2} + n+1) = \binom{n+1}{2} = \dim \bigwedge^2 k^{n+1},\]
and so on.

Now, applying $LF$ to these projective resolutions, we will get truncated Koszul complexes, which only have cohomology in the leftmost degree, which is just the kernel. So, we need the following proposition.

\begin{prop}
    Let $V=k^{n+1}$, and
    \[0 \to \bigwedge^{n+1} V \otimes_k \cO(-n-1) \xrightarrow{d_{n+1}} \cdots \to \bigwedge^2 V \otimes_k \cO(-2) \xrightarrow{d_2} \bigwedge^1 V \otimes_k \cO(-1) \xrightarrow{d_1} \bigwedge^0 V \otimes_k \cO \xrightarrow{d_0} 0\]
    the Koszul complex. Then, $\ker d_i = \Omega^i$.
\end{prop}
\begin{proof}
    Clearly $\ker d_0 = \cO = \Omega^0$, and the Euler sequence tells us $\ker d_1 = \Omega^1$.

    In general, we will still analyze the Euler sequence
    \[0 \to \Omega^1 \to V \otimes_k \cO(-1) \to \cO \to 0.\]
    In general, it is hard to say how exterior powers of vector bundles relate to each other in a short exact sequence
    \[0 \to L \to M \to N \to 0.\]
    We can, however, at least say that $\bigwedge^p M$ has a filtration
    \[\bigwedge^p M = F^0 \supseteq F^1 \supseteq \cdots \supseteq F^{p+1} = 0,\]
    where $F^i/F^{i+1} \cong \bigwedge^i L \otimes \bigwedge^{p-i} N$ by setting $F^i$ to be the subbundle of $M$ where at least $i$ wedges come from $L$.

    In our situation, all the quotients $F^i/F^{i+1}$ will be trivial except for $i=p-1,p$ because $N=\cO$ is a line bundle. The remaining interesting bit of the filtration is
    \[\bigwedge^p V\otimes_k \cO(-p) = F^{p-1} \supseteq F^p \supseteq 0,\]
    which gives us a short exact sequence
    \[0 \to \Omega^p \to \otimes \bigwedge^p V\otimes_k \cO(-1) \to \Omega^{p-1} \to 0,\]
    and indeed one can check that the composition (where we inductively assume $\Omega^{p-1}$ is the kernel of the next differential)
    \[\bigwedge^p V\otimes_k \cO(-p) \twoheadrightarrow \Omega^{p-1} \hookrightarrow \bigwedge^{p-1} V \otimes_k \cO(-p+1)\]
    is the map in the Koszul complex. Thus, we have computed the kernel.
\end{proof}

\subsection{Computation: RG on Projective Space}

The functor
\[RG = \RHom(T,-): D^b(X) \to D^b(\rmod A)\]
is harder to compute on the nose---but we can settle for computing the cohomology of $RG(\cE^\bullet)$.

In the case $X=\P^n$ with the tilting sheaf $T=\bigoplus_{i=-n}^0 \cO(i)$, and considering
\[\Coh(\P^n) \hookrightarrow D^b(\P^n) \qquad \rmod A \hookrightarrow D^b(\rmod A)\]
as degree 0 complexes, we can characterize $RG(\Coh(\P^n)) \cap \rmod A$ as follows.

\begin{prop}
    A sheaf $\cF \in \Coh(\P^n)$ has $RG(\cF)\in \rmod A$ a single quiver representation (i.e., a degree 0 complex) exactly when
    \[H^j(\cF(i)) = 0\]
    for $i=0,\dots,n$ and $j>0$.
\end{prop}

\begin{proof}
    To check $RG(\cF) \in \rmod A$, we need to show
    \[H^j(RG(\cF)) = 0\]
    for $j\neq 0$. Now, to compute $RG$, we have a have the following commutative diagram (composition of derived functors because $\cHom(T,-)$ sends injective sheaves to flasque sheaves).
    \[\begin{tikzcd}[column sep=2cm]
        D^b(\QCoh(\P^n)) \rar["R\cHom(T{,}-)"] \ar[rr,bend right,"R\Hom(T{,}-)"'] & D^b(\rMod\cEnd(T)) \rar["R\Gamma"] & D^b(\rMod A)
    \end{tikzcd}\]
    So, to compute $RG(\cF)$, we first compute
    \[R\cHom(T,\cF) = \cHom(T,\cF),\]
    since $T$ is locally free. Then, we can compute
    \begin{align*}
        R^j \Gamma(\cHom(T,\cF)) &= H^j(\cHom(\bigoplus_{i=-n}^0 \cO(i), \cF))\\
        &= H^j(\bigoplus_{i=0}^n \cF(i))\\
        &= \bigoplus_{i=0}^n H^j(\cF(i)).
    \end{align*}
    So, $RG(\cF) \in \rmod A$ exactly when each summand vanishes for $j>0$, which is statement of the proposition.
\end{proof}

\begin{rmk}
    We can get this vanishing via an assumption on the \textbf{Castelnuovo-Mumford regularity} of $\cF$. In particular, when $\cF$ is 1-regular, we get
    \[H^j(\cF(1+p-j))=0, \qquad p\geq 0,\]
    so by taking $p=i+j$ for $i=0,\dots,n$ and $j>0$, we get all the desired cohomology vanishing.
\end{rmk}


\section{Duality of Simples and Projectives}

The first thing is to establish that for an exceptional object $E\in \cD$, the inclusion $\<E\> \cong D^b(\C) \hookrightarrow \cD$ is \textit{admissible}.

\subsection{Admissibility and Mutations}

\begin{defn}
    A triangulated subcategory $\cA \subseteq \cD$ is \textbf{left (resp. right) admissible} if the inclusion
    \[i: \cA \hookrightarrow \cD\]
    admits a left (resp. right) adjoint $i^*$ (resp. $i^!$).
\end{defn}

The adjoint should be thought of as a sort of projection onto $\cA$, and the cone over the unit/counit to be the orthogonal complement.

For example, suppose $i: \cA\hookrightarrow \cD$ is left admissible, so that there is a left adjoint $i^*: \cD\to \cA$. Then for an object $X\in \cD$, we have the unit map
\[X \to i i^* X\]
comparing $X$ to the projection of $X$ onto $\cA$, which we can extend to a distinguished triangle
\[B \to X \to i i^* X \to B[1].\]
Now, observe that for any $A\in \cA$, after applying $\Hom(-, iA)$ to get a LES, we get
\[\Hom_\cD(i i^* X, i A) = \Hom_\cA(i^* X, A) = \Hom_\cD(X, i A),\]
so by the LES in $\Hom_\cD(-,iA)$, we get $\Hom_\cD(B, iA) = 0$, i.e., $B\in {}^\perp\cA$.

\begin{defn}
    Let $X\in \cD$, $i:\cA \to \cD$ the inclusion, and $i^*,i^!$ the left/right adjoints (if they exist). Denote
    \[L_\cA(X) \defeq C(ii^! X \to X) \in \cA^\perp.\]
    \[R_\cA(X) \defeq C(X \to i i^* X)[-1] \in {}^\perp\cA\]
    which are called the \textbf{left/right mutations of $X$ through $\cA$}.
\end{defn}

Maybe more intuitive notation, keeping in mind the inner product space analogies, would be something like this.
\[\mathrm{proj}^\perp_\cA X \defeq i^* X, \qquad \mathrm{orth}^\perp_\cA X \defeq L_\cA(X),\]
\[{}^\perp\mathrm{proj}_\cA X \defeq i^! X, \qquad {}^\perp\mathrm{orth}_\cA X \defeq R_\cA(X).\]

So far $i^*,i^!$, if they exist, are functorial (either by definition, or if we just define $i^*,i^!$ objectwise by satisfying a universal property from the adjunction, we also automatically get functoriality).

However, it is not immediately clear if $L_\cA, R_\cA$ are functorial---they are at least well-defined up to isomorphism because cones are well-defined up to isomorphism, but famously cones in general lack functoriality. So, we should hope something special is going on here.

\begin{prop}
    Mutations assemble into functors $L_\cA: \cD\to {}^\perp\cA$ and $R_\cA: \cD \to \cA^\perp$.
\end{prop}
\begin{proof}
    We just show $R_\cA$ is functorial. Let $f:X\to Y$ be a morphism in $\cD$, so we get the following commutative diagram by the functoriality of $i^!$ and TR3.
    \[\begin{tikzcd}
        i^!X \rar \dar["{i^* f}"] & X \rar \dar["f"] & R_\cA(X) \rar \dar[dashed] & i^* X[1] \dar["{i^! f[1]}"]\\
        i^!Y \rar & Y \rar & R_\cA(Y) \rar & i^! Y[1]\\
    \end{tikzcd}\]
    However, we do not know if there is a unique map making this diagram commute. By subtracting, we can assume $f=0$, $i^! f=0$, so our task is to show that
    \[R_\cA(X) \xrightarrow{0} R_\cA(Y)\]
    is the unique map making the diagram commute. By Yoneda, we only need to show
    \[\Hom_{{}^\perp\cA}(B, R_\cA(X)) \to \Hom_{{}^\perp\cA}(B, R_\cA(Y))\]
    is 0 for all $B\in {}^\perp\cA$. Hitting the morphism of distinguished triangles with $\Hom(B,-)$, we get
    \[\begin{tikzcd}
        0 \rar & \Hom(B,X) \rar \dar & \Hom(B,R_\cA(X)) \rar \dar & 0\\
        0 \rar & \Hom(B,Y) \rar & \Hom(B,R_\cA(Y)) \rar & 0
    \end{tikzcd}\]
    and so $\Hom(B,X)\xrightarrow{0} \Hom(B,Y)$ forces our desired morphism to be 0.

    A similar argument then shows that $L_\cA$ is functorial.
\end{proof}

One other key result about mutations.

\begin{prop}
    The restrictions $L_\cA: {}^\perp\cA \to \cA^\perp$ and $R_\cA: \cA^\perp \to {}^\perp\cA$ form an equivalence of categories.
\end{prop}

\begin{proof}
    {\color{red} write down a proof.}
\end{proof}


\begin{eg}
    Let $S$ be a smooth scheme over $k$ and $E$ a vector bundle. Consider the projective bundle $\pi: \P(E) \to S$. Then,
    \[\pi^*: D^b(S) \hookrightarrow D^b(\P(E))\]
    is full faithful, and this is a right admissible subcategory (it will also be left admissible by Grothendieck duality).

    To see this, first we check that $\pi^*$ is fully faithful, we should check
    \[\Hom(E,F) \xrightarrow{\pi^*} \Hom(\pi^* E, \pi^* F)\]
    is an isomorphism, and to check this we refer to the triangle
    \[\begin{tikzcd}
        \Hom(E,F) \rar \ar[dr] & \Hom(\pi^* E, \pi^* F) \dar\\
        & \Hom(E, R\pi_* \pi^* F)
    \end{tikzcd}\]
    where the vertical map comes from the $\pi^* \vdash R\pi_*$ adjunction, and the diagonal map is post-composition with the unit of the adjunction. So, to see the horizontal map is an isomorphism, we just check that the diagonal map is an isomorphism, so we need the unit to be an isomorphism.

    To see that $E \to R\pi_* \pi^* E$, it suffices to prove
    \[\cO_S \to R\pi_* \cO_{\P(E)}\]
    is an isomorphism, since we can just tensor it with $E$ and apply the projection formula.
    
    We work locally. Over an affine open $\Spec A$ that trivializes $\P(E)$, the map $\pi$ is the structure map
    \[\pi: \P^r_A \to \Spec A\]
    of projective space, and $R\pi_*$ is $R\Gamma$, i.e., sheaf cohomology (up to tilde). Now, this is something well-known:
    \[R\Gamma \cO_{\P(E)}|_{\Spec A} = A,\]
    which is $\cO_S|_{\Spec A}$ (up to tilde).

    So, $\pi^*$ is fully faithful, and now it is right admissible because $\pi^*$ has adjoint $R\pi_*$.
\end{eg}


\subsection{Exceptional Implies Admissible}

\begin{prop}
    For an exceptional object $E\in \cD$, the inclusion
    \[\<E\> = D^b(\mathrm{pt}) \hookrightarrow \cD\]
    is left and right admissible.
\end{prop}
\begin{proof}
    Let $X\in \cD$, we define
    \[i^! X \defeq \Hom^\bullet_\cD(E,X) \otimes E.\]
    (This seems like a reasonable candidate for projecting $X$ onto $\<E\>$, and it's functorial). To see that this is a right adjoint to $i:\<E\>\hookrightarrow \cD$, we have to produce a natural isomorphism
    \[\Hom(E,i^! X) \cong \Hom(i E,X)\]
    for $X\in \cD$. A priori, the right argument should be any object in $\<E\>$, but these are all just direct sums $\bigoplus_j E^{a_j}[j]$.

    Now, this natural isomorphism follows from the definition of $i^! X$, as
    \[\Hom(E,\Hom^\bullet(E,X)\otimes E) \cong \Hom(E,E) \otimes \Hom^\bullet(E,X) \cong \Hom^\bullet(E,X).\]
    Explicitly, this isomorphism works by interpreting a 
    \[\phi: E \to \Hom^\bullet_\cD(E,X)\otimes E,\]
    as a column vector valued in $\Hom(E,E)=k1_E$ indexed by a basis of $\Hom^\bullet_\cD(E,X)$, so a formal sum
    \[\phi = \sum_{\psi \in \Hom^\bullet(E,X)} c_\psi 1_E \otimes \psi,\]
    which is then mapped to evaluation of the formal sum
    \[\sum_{\psi \in \Hom^\bullet(E,X)} c_\psi \psi.\]
    
    So, in this way, we see $\<E\>$ is right admissible, and similarly we can define
    \[i^* X \defeq \Hom^\bullet_\cD(X,E)^* \otimes E\]
    to see $\<E\>$ is left admissible.
\end{proof}

Now, let's understand how to mutate through an exceptional object, so we need to understand the unit/counit for these adjunctions.

To study the left mutation, we look at the right adjoint
\[i^! = \Hom^\bullet(E,-)\otimes E.\]
The counit is given by the image of $1_{i^! X}$
\[\Hom(i^! X, i^! X) \xrightarrow{\cong} \Hom(i i^! X, X).\]
If we consider
\[1_{i^! X} \in \Hom(i^! X, i^! X) = \Hom(\Hom^\bullet(E,X) \otimes E, \Hom^\bullet(E,X) \otimes E)\]
as the identity matrix, in the sense of a matrix valued in $\Hom(E,E)$ with rows and columns indexed by a basis $\Hom^\bullet(E,X)$, then $1_{i^! X}$ is sent to ``summing the columns'', which is nothing other than evaluation
\[\Hom^\bullet(E,X) \otimes E \longrightarrow X.\]
So, $L_E(X)$ is defined by the distinguished triangle
\[\Hom^\bullet(E,X) \otimes E \xrightarrow{\eval} X \to L_E(X).\]

To study the right mutation, we look at the left adjoint
\[i^* = \Hom^\bullet(-,E)^* \otimes E.\]
The unit is then given by the image of $1_{1^* X}$
\[\Hom(i^* X, i^* X) \xrightarrow{\cong} \Hom(X, ii^* X).\]
The identity morphism on $\Hom^\bullet(X,E)^* \otimes E$, after picking a basis
\[\Hom^\bullet(X,E)=\<\psi_1,\dots,\psi_n\>\]
is the identity morphism on each $\psi_i^* \otimes E$ component. Then, the adjoint morphism
\[X \xrightarrow{\mathrm{coeval}} \Hom^\bullet(X,E)^* \otimes E\]
is defined on each component by the map $X\to \psi_i^* \otimes E$ given by $\psi_i$. Interpreted coordinate free, this is ``coevaluation'', which is the curried evaluation map
\[X \otimes \Hom^\bullet(X,E) \to E.\]

\subsection{Dual Exceptional Collections}




\nocite{*} % Include all references in refs.bib regardless of whether they are referenced or not
\bibliographystyle{plain} % We choose the "plain" reference style
\bibliography{tilting_refs} % Entries are in the refs.bib file

\end{document}
























% \section{Beilinson's Argument}

% The main result we are working towards is the following.
% \begin{thm}
%     The following collections of sheaves on $\P^n$ are strong full exceptional sequences:
%     \[\cO(-n),\cO(-n+1),\dots,\cO\]
%     and
%     \[\cO,\Omega^1(1),\dots,\Omega^n(n).\]
% \end{thm}

% We will also be able to hit this collection with any auto-equivalence of $D^b(\P^n)$ to get more exceptional collections. To see this, we are going to lightly dabble in Fourier-Mukai transforms, but nothing too fancy---we'll ``just'' study the identity functor $1_{D^b(\P^n)}$.

% \begin{defn}
%     Let $X,Y$ be smooth projective varieties, and denote the projections
%     \[\begin{tikzcd}[column sep=0cm]
%         & X\times Y \ar[dl,"q"'] \ar[dr,"p"] &\\
%         X & & Y
%     \end{tikzcd}\]
%     Let $\cP$ be a complex of sheaves in $D^b(X\times Y)$, called a \textbf{Fourier-Mukai kernel}, and define the corresponding \textbf{Fourier-Mukai transform}
%     \begin{align*}
%         \Phi_{\cP}: D^b(X) &\longrightarrow D^b(Y)\\
%         \cE &\longmapsto Rp_*(Lq^*\cE \otimes^L \cP).
%     \end{align*}
% \end{defn}

% \begin{prop}
%     Let $\cO_{\Delta}$ be the structure sheaf of the diagonal on $X\times X$, i.e., if $\Delta: X\to X\times_k X$ is the diagonal, $\cO_{\Delta} = \Delta_* \cO_X$. Then,
%     \[\Phi_{\cO_\Delta} \cong 1_{D^b(X)}.\]
% \end{prop}

% \begin{proof}
%     First, we note that $\Delta_*\cO_X$, considered as a degree 0 complex, coincides with $R\Delta_* \cO_X$ where $\cO_X$ is considered as a degree 0 complex, since $\Delta$ is a closed immersion (so an affine map).
    
%     We then compute, by the derived projection formula (which is true even for non-locally free sheaves, since every complex of coherent sheaves has a locally free replacement) the following:
%     \begin{align*}
%         \Phi_{\cO_\Delta}(\cE) &= Rp_*(Lq^*\cE \otimes^L \Delta_* \cO_X)\\
%         &= Rp_*(R\Delta_*(\cO_X \otimes^L L\Delta^*Lq^* \cE))\\
%         &= R(p \circ \Delta)_*(\cO_X \otimes^L L(q\circ \Delta)^* \cE)\\
%         &= R(1_X)_*(\cO_X \otimes^L L(1_X)^* \cE)\\
%         &= \cE
%     \end{align*}
%     (where all the equalities are really only natural isomorphisms). So, indeed $\Phi_{\cO_\Delta} \cong 1_{D^b(X)}$.
% \end{proof}

% So far, we've come up with a worse way to write the identity functor. We'll make it even worse by replacing $\cO_\Delta$ with a quasi-isomorphic complex, which will be a locally free resolution. This is commonly referred to as \textbf{Beilinson's resolution of the diagonal}.

% \begin{prop}
%     Write $\P = \P^n$. Define
%     \[\cV = \cO(-1) \boxtimes \Omega^1(1)\]
%     on $\P\times \P$. There is a finite locally free resolution of the sheaf $\cO_{\Delta}$ on $\P\times \P$ given by
%     \[0 \to \bigwedge^n \cV \to \cdots \to \bigwedge^1 \cV \to \bigwedge^0 \cV \to \cO_\Delta \to 0.\]
% \end{prop}

% \begin{proof}
%     We will produce a section
%     \[\eval \in \Gamma(\mathcal{V}^\vee),\]
%     such that the vanishing closed subscheme of $\eval$ is $\Delta$. This implies that
%     \[\coker(\eval: \cV\to \cO_{\P\times \P}) \cong \cO_\Delta,\]
%     and because
%     \[\dim \Delta = \mathrm{rank} \Omega^1(1) = \mathrm{rank} \cV,\]
%     this implies that the corresponding Koszul complex
%     \[0 \to \bigwedge^n \cV \to \cdots \to \bigwedge^1 \cV \to \bigwedge^0 \cV \to \cO_\Delta \to 0\]
%     is exact.

%     So, the crux is defining $\eval$. We can start by defining $\eval$ on a larger sheaf,
%     \[\cO_\P^{n+1} \boxtimes (\cO_\P^{n+1})^\vee \to \cO_{\P\times \P}\]
%     and then embed $\cO(-1) \hookrightarrow \cO^{n+1}$ as well as $\Omega^1(1) \hookrightarrow (\cO^{n+1})^\vee$. We'll show that this restriction is exactly the $\eval$ we desire.

%     We get $\cO(-1)\hookrightarrow \cO^{n+1}$ by tensoring
%     \[\cO \xrightarrow{\begin{bmatrix}
%         x_0\\\vdots\\x_n
%     \end{bmatrix}} \cO(1)^{n+1}\]
%     by $\cO(-1)$
% \end{proof}

% We can then compute
% \[\bigwedge^i \cO(-1) \boxtimes \Omega^1(1) = \cO(-i) \boxtimes \Omega^i(i)\]




% \newpage
% \section{Appendix}

% Here we prove some standard algebraic geometry facts, just to collect them.

% \subsection{Koszul Complex of Vector Bundle}

% \begin{prop}
%     Let $\cE$ be a locally free sheaf of rank $r$ on a scheme $X$, and $s\in \Gamma(\cE^\vee)$ a global section, which gives a closed subscheme $V(s)$. Then,
%     \[\coker(\cE \xrightarrow{s} \cO_X) \cong \cO_{V(s)}.\]
% \end{prop}

% \begin{proof}
%     As usual,
%     \[\coker(s) = \cO_X/\im s,\]
%     and $\cO_{V(s)} = \cO_X/\cI_{V(s)}$, so we just have to show that $\im s = \cI_{V(s)}$.

%     We compare $\cI_{V(s)}$ and $\im s$ on the base of trivializing affine opens. In this way, without loss of generality, we just assume $X=\Spec A$ and $\cE=\cO_X^r$ is trivial.

%     Now, $s \in \Hom(\cO_X^r, \cO_X)$ is a choice of functions $f_1,\dots,f_r$ for the images of sections $e_1,\dots,e_r \in \Gamma(\cO_X^r)$.

%     We can then compute $\im(s) = \tilde{(f_1,\dots,f_r)}$, and the local definition of $V(s)$ is exactly as the closed subscheme $\Spec A/(f_1,\dots,f_r)$, which again corresponds to the ideal sheaf $\tilde{(f_1,\dots,f_r)}$.
% \end{proof}

% \begin{prop}
%     Let $\cE$ be a locally free sheaf of rank $r$ on a scheme $X$, and $s\in \Gamma(\cE^\vee)$ a global section. Then, there is a complex called the \textbf{Koszul complex} 
%     \[0 \to \bigwedge^r \cE \to \cdots \to \bigwedge^1 \cE \to \bigwedge^0 \cE \to \cO_{V(s)} \to 0,\]
%     and if $X$ is locally Cohen-Macaulay and $\codim V(s) = r$, then it is exact.
% \end{prop}

% \begin{proof}
%     We can identify $\bigwedge^i E$ as a subsheaf of a tensor product, (as in the differential geometry-esque style of wedge product as alternating functionals on the dual)
%     \[\bigwedge^i E \hookrightarrow \bigwedge^{i-1} E \otimes E,\]
%     by hitting the following morphism of presheaves with sheafification
%     \[e_1 \wedge \cdots \wedge e_i \longmapsto \sum_{\ell=1}^i (-1)^\ell e_1 \wedge \cdots \wedge \hat{e_\ell} \wedge \cdots \wedge e_i \otimes e_\ell.\]
%     Then, we have a morphism
%     \[\bigwedge^{i-1} E \otimes E \to \bigwedge^{i-1} E \otimes \cO_X \cong \bigwedge^{i-1} E\]
%     from $E\xrightarrow{s} \cO_X$. Now, the composition of these morphisms is
%     \[\bigwedge^i E \to \bigwedge^{i-1} E\]
%     which is referred to as contraction with $s$.

%     Now, on stalks, choosing an identification $E_x \cong \cO_{X,x}^r$, our morphism $s: E_x \to \cO_{X,x}$ gets identified with a morphism $\cO_{X,x}^r \to \cO_{X,x}$ defined by sending the basis $e_1,\dots,e_r$ to germs $f_1,\dots,f_r$. Then, under this identification, the complex we defined is exactly the usual Koszul complex $K(f_1,\dots,f_r)$ on the local ring $\cO_{X,x}$.

%     So, this tells us in particular that our sequence of morphisms between sheaves $\bigwedge^i E$ is a complex. We can also say something about exactness. Assume now that $X$ is locally Cohen-Macaulay and $\codim V(s)=r$, and let $x\in X$ be a closed point. Then, we want to know that $f_1,\dots,f_r$ is a regular sequence, but this follows from \cite[Theorem 2.1.2]{Bruns_Herzog_1998}, since $\dim \cO_{X,x}/(f_1,\dots,f_r) = r$ by our codimension assumption. So, our complex is exact when we take stalks at any closed point, so our complex is exact.
% \end{proof}

% \subsection{Tautological Bundle}

% \subsubsection{Geometric Vector Bundle Computation}
% The tautological bundle on $\P^n$, classically, is the \textit{geometric} vector bundle defined as a subbundle of
% \[\pi: \P^n \times \A^{n+1} \longrightarrow \P^n\]
% which consists of pairs $([x],v)$ such that $v\sim x$, i.e., considering $\P^n$ as the set of 1-dimensional subspaces $\ell$ in $\A^{n+1}$, we are setting the fiber over $\ell$ to be $\ell$ itself.

% To define this in a more scheme theoretic way, we take the line bundle
% \[\cO_{\P^n}(1) \boxtimes \cO_{\A^{n+1}}\]
% (where we give $\P^n$ homogeneous coordinates $x_0,\dots,x_n$ and $\A^{n+1}$ coordinates $v_0,\dots,v_n$), and we look at the vanishing locus of the sections
% \[x_i v_j - x_j v_i, \qquad i,j=0,\dots,n,\]
% which is exactly asserting that $v$ is on the line spanned by $x$.

% Denote
% \[L = V(x_i v_j - x_j v_i)\]
% the zero scheme of the sections. We can check using the affine open cover $D_+(x_i) \times \A^{n+1}$ that $L$ is integral. We then restrict $\pi$ to
% \[p: L \to \P^n,\]
% and we claim that this is a geometric vector bundle, and that its sheaf of sections is exactly $\cO(-1)$. As a result, we'll also see how to embed $\cO_{\P^n}(-1) \hookrightarrow \cO_{\P^n}^{n+1}$.

% We can trivialize this vector bundle over $D_+(x_i)\subseteq \P^n$, since it has preimage $p^{-1}(D_+(x_i))$ which is the vanishing closed subscheme in $D_+(x_i)\times \A^{n+1}$ of the restricted sections of the line bundle, so we get
% \[p^{-1}(D_+(x_i)) = \Spec k\left[\frac{x_0}{x_i},\dots,\frac{x_n}{x_i}\right][v_0,\dots,v_n]/(v_j - \frac{x_j}{x_i} v_i) \cong \Spec k\left[\frac{x_0}{x_i},\dots,\frac{x_n}{x_i}\right][v_i],\]
% which is exactly a copy of $\A_{D_+(x_i)}^1$, and this isomorphism respects the projections in the sense that the following diagram commutes.
% \[\begin{tikzcd}
%     p^{-1}(D_+(x_i)) \dar \rar["f_i"] & \A_{D_+(x_i)}^1 \ar[dl]\\
%     D_+(x_i) &
% \end{tikzcd}\]
% One then computes that
% \[\Spec \cO_{\P^n}(D_+(x_i x_j))[v] \cong \A_{D_+(x_i x_j)}^1 \xrightarrow{f_j^{-1} \circ f_i} \A_{D_+(x_i x_j)}^1 \cong \Spec \cO_{\P^n}(D_+(x_i x_j))[v]\]
% is given by the map on rings
% \[\cO_{\P^n}(D_+(x_i x_j))[v] \to \cO_{\P^n}(D_+(x_i x_j))[v]\]
% sending $v\mapsto \frac{x_j}{x_i} v$, which is $D_+(x_i x_j)$-linear. Thus, $p: L\to \P^n$ is a geometric vector bundle.

% Now, we can get the corresponding invertible sheaf by gluing along these same transition maps. Recall that the sheaf of sections of the projection $\A^1_{D_+(x_i)} \to D_+(x_i)$ gives $\cO_{D_+(x_i)}$, since along an affine open $\Spec A$ (these form a basis of open sets) a section of this projection is a morphism of $A$-schemes
% \[\Spec A \xrightarrow{s} \Spec A[v]\]
% which is exactly given by picking an image $s\in A$ for $v$, so we recover sections over $\Spec A$ as $A = \cO_{D_+(x_i)}(\Spec A)$.

% Unraveling this correspondence, the sections of this projection carry the $\cO_{\P^n}$ multiplication structure as follows. On an affine open $\Spec A$, with $f\in \cO_{\P^n}(\Spec A)=A$, we send a section $\Spec A \xrightarrow{s} \Spec A[v]$ (corresponding to $v\mapsto s$) to the section $v\mapsto f\cdot s$.

% Now, we glue these trivial bundles together by pushing forward sections, i.e., if we have a section
% \[D_+(x_i x_j) \to \A^1_{D_+(x_i x_j)},\]
% we can post compose it with $f_j \circ f_i^{-1}$ to get another section, giving a morphism
% \[\cO_{D_+(x_i x_j)} \xrightarrow{\tau_{ij}} \cO_{D_+(x_i x_j)}.\]
% Now, to compose we look at maps on coordinate rings, i.e., to compute the composition, for a section $s$,
% \[\Spec \cO_{\P^n}(D_+(x_i x_j)) \xrightarrow{s} \Spec \cO_{\P^n}(D_+(x_i x_j))[v] \xrightarrow{f_j \circ f_i^{-1}} \Spec \cO_{\P^n}(D_+(x_i x_j))[v],\]
% we go backwards
% \[\cO_{\P^n}(D_+(x_i x_j))[v] \xrightarrow{\Gamma(f_j\circ f_i^{-1})} \cO_{\P^n}(D_+(x_i x_j))[v] \xrightarrow{\Gamma(s)} \cO_{\P^n}(D_+(x_i x_j)),\]
% where we computed $f_j \circ f_i^{-1}$ sends $v\mapsto \frac{x_j}{x_i}$, which is subsequently mapped to $\frac{x_j}{x_i}s$. So, the morphism $\tau_{ij}$ is exactly multiplication by $\frac{x_j}{x_i}$.

% Now, the gluing data $(D_+(x_i), \tau_{ij} = \frac{x_j}{x_i})$ is exactly that of $\cO(-1)$, so indeed $\cO(-1)$ is the rightful heir to the title tautological bundle.

% \subsubsection{Embed in Trivial Vector Bundle}

% By construction, we had a commutative diagram
% \[\begin{tikzcd}
%     L \ar[dr] \rar & \P^n \times \A^{n+1} \dar\\
%     & \P^n
% \end{tikzcd}\]
% i.e., a morphism of vector bundles, which in turn gives rise to a morphism of sections by post composing a section $\P^n \to L$ with the morphism $L\to \P^n \times \A^{n+1}$.

% We unravel this definition to see exactly how this is embedding $\cO(-1) \hookrightarrow \cO^{n+1}$ on the trivializing open cover $\{D_+(x_i)\}$.

% Restricted to $D_+(x_i)$, we get
% \[\begin{tikzcd}
%     \Spec \cO_{\P^n}(D_+(x_i))[v] \rar & \Spec \cO_{\P^n}(D_+(x_i))[v_0,\dots,v_n]
% \end{tikzcd}\]
% where the map on rings sends $v_j \mapsto \frac{x_j}{x_i} v$ (where when $j=i$, $v_i\mapsto v$). This corresponds to
% \begin{align*}
%     \cO_{D_+(x_i)} &\longrightarrow \cO_{D_+(x_i)}^{n+1}\\
%     1 &\longmapsto (\frac{x_0}{x_i}v, \dots, \frac{x_n}{x_i}v),
% \end{align*}
% and realizing $\cO(-1)|_{D_+(x_i)} = \frac{1}{x_i}\cO_{D_+(x_i)}$, the map on $\cO(-1)$ is given by
% \begin{align*}
%     \cO(-1)|_{D_+(x_i)} &\longrightarrow \cO_{D_+(x_i)}^{n+1}\\
%     f &\longmapsto (x_0 f, \dots, x_n f),
% \end{align*}
% so in other words, this is exactly
% \[\cO \xrightarrow{\begin{bmatrix}
%     x_0\\\vdots\\x_n
% \end{bmatrix}} \cO(1)^{n+1}\]
% twisted by $-1$.

























% \subsection{(Geometric) Vector Bundles}

% \begin{defn}
%     Let $X$ be a scheme. A \textbf{pre-vector bundle} on $X$ is an $X$-scheme $E\xrightarrow{p} X$ such that for every $X$-scheme $T$, the set of $T$-valued points $E(T)$ carries the structure of a module over $\cO_T(T)$ (this is included in the data of a pre-vector bundle). We require that for $X$-morphisms $T\to T'$, the induced map (precomposition) $E(T')\to E(T)$ is linear over $\cO_{T'}(T') \to \cO_T(T)$.
% \end{defn}

% In particular, when we take the $X$-scheme which is an open $U\subseteq X$, we are requiring that sections of $p^{-1}(U) \to U$ can be scaled by functions on $U$

% This section is essentially just a solution to Hartshorne exercise II.5.18.

% \begin{defn}
%     A \textbf{(geometric) vector bundle} of rank $r$ over a scheme $X$ is a scheme $E$ and a morphism $p: E\to X$ together with the data of an open cover $\{U_i\}$ of $X$ together with \textbf{trivializations}, i.e., isomorphisms $\phi_i$ fitting into commutative diagrams
%     \[\begin{tikzcd}
%         \pi^{-1}(U_i) \ar[d,"p"] \rar["\phi_i"] & \A_{U_i}^r \ar[dl]\\
%         U_i & 
%     \end{tikzcd}\]
%     such that the compositions
%     \[\begin{tikzcd}
%         \A_{U_i\cap U_j}^r \rar["\phi_i^{-1}"] & p^{-1}(U_i\cap U_j) \rar["\phi_j"] & \A_{U_i \cap U_j}^r
%     \end{tikzcd}\]
%     are linear maps, i.e., for any $\Spec A \subseteq U_i\cap U_j$ affine open, we get
%     \[\A_A^r \to \A_A^r\]
%     induced by a map of rings
%     \begin{align*}
%         A[x_1,\dots,x_r] &\longrightarrow A[x_1,\dots,x_r]\\
%         x_j &\longmapsto \sum_{i=1}^r a_{ij} x_i
%     \end{align*}
%     for some matrix $(a_{ij})$ in $A$.
% \end{defn}

% It is slightly annoying that we have to include the data of the trivializations in the definition of a vector bundle, but this is our way of remembering the linear structure on the fibers.

% \begin{defn}
%     Let $(E,p,\{U_i,\phi\})$ and $(F,q,\{W_j,\psi_j\})$ be vector bundles on $X$. A \textbf{morphism of vector bundles} is a morphism $F:E\to F$ commuting with projections, i.e., $q\circ F = p$, such that
% \end{defn}
















% \begin{prop}
%     Suppose $E$ classically generates a triangulated category $\cD$, and $F: \cD\to \cD$ is an exact functor which sends $E$ to an object isomorphic to $E$ and is identity on $\Hom_{\cD}(E,E[i]) \to \Hom_{\cD}(E,E[i])$ (after pre/post composing with the isomorphism).

%     Then, $F$ is naturally isomorphic to $1_\cD$.
% \end{prop}
% \begin{proof}
%     We will use induction as in Remark \ref{induction-triangulated-cat}. Let $T$ be the property of objects $A\in \cD$ that $F: \Hom(E,A) \to \Hom(FE,FA)$ is an isomorphism, and we check properties (1)-(4).

%     The base case (1) holds by assumption, (2) and (3) hold because $T$ is additive, and (4) follows from the 5-lemma, where we get a morphism of LES after applying $\Hom(E,-)$ and $\Hom(FE,-)$ to the distinguished triangle and its image under $F$. Thus, $T$ holds for all of $\cD$.

%     Now, fixing $B$ for the right argument of $\Hom$, same strategy proves show for the property $T_B$ for all of $\cD$, where $T_B(A)$ iff $\Hom(B,A) \to \Hom(FB,FA)$ is an isomorphism---so $F$ is fully faithful.

%     Next, we show that $A\cong FA$ in such a way that for any morphism $E\to A$ we have the commutative diagram
%     \begin{cd}
%         E \dar\rar & A\dar \\
%         F(E) \rar & F(A)
%     \end{cd}
%     (which is only part of what we will need to assemble $A\to FA$ into a natural transformation).
    
%     Denote this property by $S$. The base case (1) holds by assumption, (2) and (3) because $T$ is additive, and we check (4) on a distinguished triangle $K\to L \to M \to L[1]$. For example, assume $P(K), P(L)$, so by TR3 for triangulated categories we get
%     \begin{cd}
%         K \rar \dar & L \rar \dar & M \rar \dar[dashed] & K[1] \dar\\
%         FK \rar & FL \rar & FM \rar & FK
%     \end{cd}
%     so we at least get a morphism $M\to FM$, and this is an isomorphism by hitting the diagram with the Yoneda embedding and applying the 5-lemma. We now need to check that this commutes with any morphism $E\to M$??

%     Now, we want to assemble these isomorphisms (or variations of them) into a natural transformation $1_\cD \Rightarrow F$. To do this, we check
    
    
%     we fix an isomorphism $\phi_A: A\to FA$ for every object $A\in \cD$, and we build our natural transformation.
    
%     Now, we try to build a natural transformation $\eta: 1_\cD \to F$ by finding components $\eta_A: A \to FA$. To find these components, we study the composition
%     \[\Hom(E,A) \xrightarrow{F} \Hom(FE,FA) \to \Hom(E,FA).\]
%     Denote $P(A)$ the property that the composition is induced by pushforward along some map $\eta_A: A\to FA$. We again prove $P$ by a sort of induction.

%     To show the base case (1), we see for $A=E[i]$, we have
%     \[\Hom(E,E[i]) \xrightarrow{F} \Hom(FE,FE[i]) \to \Hom(E,FE[i])\]
%     which post composes to be the identity $\Hom(E,E[i])$, so pulling back this must be post composition by our isomorphism $E[i]\to FE[i]$.

%     Next, (2) and (3) follow because everything here is additive.
%     % \begin{cd}
%     %     \Hom(E{,}K\oplus L) \dar["(\pi_1)_*{,}(\pi_2)_*"] \rar["F"] & \Hom(E{,}F(K\oplus L)) \dar["(\pi_1)_*{,}(\pi_2)_*"]\\
%     %     \Hom(E{,}K)\oplus \Hom(E{,}L) \rar["F\oplus F"] & \Hom(E{,}F(K)) \oplus \Hom(E{,}F(L))
%     % \end{cd}

%     Finally to show (4), let $K\to L \to M \to K[1]$ be a distinguished triangle, and assume $P(K),P(M)$. Now, build a diagram
%     \begin{cd}
%         \cdots \rar & \Hom(E,K) \rar \dar["(\eta_K)_*"] & \Hom(E,L) \rar \dar["?"] & \Hom(E,M) \rar \dar["(\eta_M)_*"] & \cdots\\
%         \cdots \rar & \Hom(E,F(K)) \rar & \Hom(E,F(L)) \rar & \Hom(E,F(M)) \rar & \cdots
%     \end{cd}
%     and then 

%     ...or should we show
    
%     is induced by pushforward along some $\eta_A: A\to FA$. Denote the property that the map is induced

%     {\color{red} I really need to show that this is natural...} Next, we show $FA \cong A$ via Yoneda, so we define the property $Y$ by saying $Y(A)$ holds iff
%     \[\Hom(-,A) \cong \Hom(-,FA).\]
%     The base case (1) holds for $Y$, since
%     \[\Hom(E,A) \cong \Hom(FE,FA) \cong \Hom(E,FA)\]
%     by pulling back along the isomorphism $FE\to E$. Next, (2) and (3) hold because $\Hom$ interacts well with biproducts. E.g.,
%     \[\Hom(K\oplus L, A) \cong \Hom(K\oplus L, FA)\]
% \end{proof}